%--------------------------------------------------------------------%
%
% Driver khusus konversi ke DOCX (pandoc).
%
% JANGAN dipakai untuk membangun PDF. Berkas ini sengaja TIDAK memuat
% if-itb-thesis.sty, daftar isi/tabel/gambar, atau \printbibliography
% karena hal-hal tersebut tidak dikenali pandoc dan ditangani oleh
% opsi pandoc (--toc, --citeproc) serta reference.docx.
%
% Halaman tata letak (sampul, logo, halaman judul, persetujuan,
% pengesahan, pernyataan) SENGAJA tidak disertakan karena hanya berupa
% tata letak LaTeX murni yang tidak dapat direproduksi di DOCX. Halaman
% tersebut tetap ada di PDF (thesis.tex). DOCX berfokus pada isi:
% Ringkasan, Summary, Kata Pengantar, bab, daftar rujukan, lampiran,
% dan riwayat hidup.
%
% PENTING: Jaga agar daftar \input di bawah tetap sinkron dengan
% src/thesis.tex. Bila menambah/menghapus bab, ubah keduanya.
%
%--------------------------------------------------------------------%

\documentclass[12pt,a4paper,onecolumn,oneside,final]{report}

%-------------------------------------------------------------------%
%
% Informasi umum skripsi — SUNTING DI SINI SAJA.
% Semua halaman front matter mengambil nilai dari makro di bawah.
%
%-------------------------------------------------------------------%

\newcommand{\judul}{Judul Skripsi yang Informatif dan Mencerminkan Isi}
\newcommand{\penulis}{Nama Lengkap Penulis}
\newcommand{\nim}{123456789}

\newcommand{\departemen}{Teknik Elektro dan Informatika}
\newcommand{\fakultas}{Fakultas Teknik}
\newcommand{\programstudi}{Program Studi Teknik Informatika}
% Diisi pada halaman judul: "...menyelesaikan program <gelarprogram>"
\newcommand{\gelarprogram}{Sarjana Teknik Informatika}

\newcommand{\bulantahun}{Juli 2026}
\newcommand{\kotatanggal}{Malang, 1 Juli 2026}

\newcommand{\pembimbingSatu}{Nama dan Gelar Pembimbing I}
\newcommand{\nipSatu}{NIP 123456789}
\newcommand{\pembimbingDua}{Nama dan Gelar Pembimbing II}
\newcommand{\nipDua}{NIP 123456789}

% Kata kunci ringkasan (3--5 kata), dipisah koma
\newcommand{\katakunci}{kata kunci satu, kata kunci dua, kata kunci tiga}
\newcommand{\katakunciEn}{keyword one, keyword two, keyword three}


\begin{document}

    %------------------------------------------------------------%
    % Bagian awal (isi; halaman tata letak sengaja dilewati)
    %------------------------------------------------------------%
    %--------------------------------------------------------------------%
%
% Ringkasan (Pedoman UM Lampiran 7a) — Bahasa Indonesia, spasi tunggal,
% maksimal 750 kata. Baris identitas: Nama-akhir, Nama-awal. Tahun. dst.
%
%--------------------------------------------------------------------%

\clearpage
\addcontentsline{toc}{chapter}{Ringkasan}

\begin{center}
    \bfseries RINGKASAN
\end{center}

\begin{singlespace}

\noindent \penulis. \bulantahun. \textit{\judul}. Skripsi, \departemen, \fakultas, Universitas Negeri Malang. Pembimbing: (I) \pembimbingSatu; (II) \pembimbingDua.

\vspace{0.5\baselineskip}
\noindent \textbf{Kata Kunci:} \katakunci%

\vspace{\baselineskip}

Dalam pemrosesan bahasa alami (NLP) model bahasa pra-latih seperti IndoBERT unggul pada berbagai tugas pemrosesan bahasa alami berbahasa Indonesia, tetapi ukuran parameternya yang besar membatasi penerapan pada perangkat dengan sumber daya terbatas. \textit{Post-Training Quantization} (PTQ) INT4 menjadi salah satu solusi kompresi, namun belum ada studi yang mengevaluasinya secara sistematis pada IndoBERT untuk domain finansial Indonesia. Penelitian ini mengevaluasi PTQ INT4 pada IndoBERT (\texttt{indobert\allowbreak\_base\allowbreak\_v1}) untuk analisis sentimen dan klasifikasi topik finansial menggunakan tiga dataset, yaitu IndoFinSent (2.274 entri), \textit{Indonesian Financial PhraseBank} (2.264 entri), dan dataset topik finansial Twitter (21.105 dokumen dengan 20 topik). Model di-\textit{fine-tune} pada presisi FP32 dengan \textit{learning rate} 3e-5, \textit{batch size} 16, dan 3 \textit{epoch}, kemudian dikuantisasi pada lapisan linear dengan cakupan \textit{weight-only} (W4A16) dan dibandingkan pada empat tingkat presisi, yaitu FP32, FP16, INT8, dan INT4. Performa diukur dengan Macro F1-Score dan accuracy, sedangkan efisiensi diukur melalui ukuran model, \textit{latency}, \textit{throughput}, dan rasio kompresi. Hasil penelitian menunjukkan bahwa [hasil eksperimen]. Temuan ini diharapkan memberikan gambaran mengenai kelayakan PTQ INT4 untuk \textit{deployment} IndoBERT pada perangkat dengan sumber daya terbatas di domain finansial Indonesia.

\end{singlespace}
\clearpage
   % Ringkasan
    %--------------------------------------------------------------------%
%
% Summary (Pedoman UM Lampiran 7 — versi Bahasa Inggris dari Ringkasan)
% Spasi tunggal, maksimal 750 kata.
%
%--------------------------------------------------------------------%

\clearpage
\addcontentsline{toc}{chapter}{Summary}

\begin{center}
    \bfseries SUMMARY
\end{center}

\begin{singlespace}

\noindent \penulis. \bulantahun. \textit{English Translation of the Thesis Title}. Thesis, \departemen, \fakultas, State University of Malang. Advisors: (I) \pembimbingSatu; (II) \pembimbingDua.

\vspace{0.5\baselineskip}
\noindent \textbf{Keywords:} \katakunciEn%

\vspace{\baselineskip}
\blindtext%

\end{singlespace}
\clearpage
   % Summary
    %--------------------------------------------------------------------%
%
% Kata pengantar
%
% @author Petra Novandi
%
%--------------------------------------------------------------------%
%
% Tuliskan ucapan terima kasih kepada pembimbing, keluarga, dan pihak
% lain yang membantu. Ganti teks placeholder di bawah dengan isi
% sebenarnya.
%
%--------------------------------------------------------------------%

\chapter*{Kata Pengantar}
\addcontentsline{toc}{chapter}{Kata Pengantar}

Gunakan bagian ini untuk memberikan ucapan terima kasih kepada semua pihak yang secara 
langsung atau tidak langsung membantu penyelesaian tugas akhir, termasuk pemberi beasiswa jika ada. 
Utamakan untuk memberikan ucapan terima kasih kepada tim pembimbing tugas akhir dan staf pengajar atau pihak program studi, bahkan sebelum mengucapkan terima kasih kepada keluarga. Ucapan terima kasih sebaiknya bukan hanya menyebutkan nama orang saja, tetapi juga memberikan penjelasan bagaimana bentuk bantuan/dukungan yang diberikan. Gunakan bahasa yang baik dan sopan serta memberikan kesan yang enak untuk dibaca. Sebagai contoh: ``Tidak lupa saya ucapkan terima kasih kepada teman dekat saya, Tito, yang sejak satu tahun terakhir ini selalu memberikan semangat dan mengingatkan saya apabila lengah dalam mengerjakan Tugas Akhir ini. Tito juga banyak membantu mengoreksi format dan layout tulisan. Apresiasi saya sampaikan kepada pemberi beasiswa, Yayasan Beasiswa, yang telah memberikan bantuan dana kuliah dan biaya hidup selama dua tahun. Bantuan dana tersebut sangat membantu saya untuk dapat lebih fokus dalam menyelesaikan pendidikan saya.\ldots' Ucapan permintaan maaf karena kekurangsempurnaan hasil Tugas Akhir tidak perlu ditulis.

\blindtext%
     % Kata Pengantar

    %------------------------------------------------------------%
    % Bagian inti
    %------------------------------------------------------------%
    %--------------------------------------------------------------------%
%
% Bab I — Pendahuluan
%
% @author Petra Novandi
%
%--------------------------------------------------------------------%
%
% Berisi latar belakang, rumusan masalah, tujuan, batasan, metodologi,
% dan sistematika pembahasan. Ganti teks placeholder dengan isi
% sebenarnya.
%
%--------------------------------------------------------------------%

\chapter{Pendahuluan}

Bab Pendahuluan secara umum yang dijadikan landasan kerja dan arah kerja penulis tugas akhir, berfungsi mengantar pembaca untuk membaca laporan tugas akhir secara keseluruhan.

\section{Latar Belakang}

Latar Belakang berisi dasar pemikiran, kebutuhan atau alasan yang menjadi ide dari topik tugas akhir. Tujuan utamanya adalah untuk memberikan informasi secukupnya kepada pembaca agar memahami topik yang akan dibahas.  Saat menuliskan bagian ini, posisikan anda sebagai pembaca – apakah anda tertarik untuk terus membaca?

\blindtext%

\section{Rumusan Masalah}

Rumusan Masalah berisi masalah utama yang dibahas dalam tugas akhir. Rumusan masalah yang baik memiliki struktur sebagai berikut:

\begin{enumerate}
    \item Penjelasan ringkas tentang kondisi/situasi yang ada sekarang terkait dengan topik utama yang dibahas Tugas Akhir.
    \item Pokok persoalan dari kondisi/situasi yang ada, dapat dilihat dari kelemahan atau kekurangannya. Bagian ini merupakan inti dari rumusan masalah.
    \item Elaborasi lebih lanjut yang menekankan pentingnya untuk menyelesaikan pokok persoalan tersebut.
    \item Usulan singkat terkait dengan solusi yang ditawarkan untuk menyelesaikan persoalan.
\end{enumerate}

Penting untuk diperhatikan bahwa persoalan yang dideskripsikan pada subbab ini akan dipertanggungjawabkan di bab Evaluasi apakah terselesaikan atau tidak.

\blindtext%

\section{Tujuan}

Tuliskan tujuan utama dan/atau tujuan detil yang akan dicapai dalam pelaksanaan tugas akhir. Fokuskan pada hasil akhir yang ingin diperoleh setelah tugas akhir diselesaikan, terkait dengan penyelesaian persoalan pada rumusan masalah. Penting untuk diperhatikan bahwa tujuan yang dideskripsikan pada subbab ini akan dipertanggungjawabkan di akhir pelaksanaan tugas akhir apakah tercapai atau tidak.

\blindtext%

\section{Batasan Masalah}

Tuliskan batasan-batasan yang diambil dalam pelaksanaan tugas akhir. Batasan ini dapat dihindari (tidak perlu ada) jika topik/judul tugas akhir dibuat cukup spesifik.

\blindtext%

\section{Metodologi}

Tuliskan semua tahapan yang akan dilalui selama pelaksanaan tugas akhir. Tahapan ini spesifik untuk menyelesaikan persoalan tugas akhir. Tahapan studi literatur tidak perlu dituliskan karena ini adalah pekerjaan yang harus Anda lakukan selama proses pelaksanaan tugas akhir.

\blindtext%

\section{Sistematika Pembahasan}

Subbab ini berisi penjelasan ringkas isi per bab. Penjelasan ditulis satu paragraf per bab buku.

\blindtext%

    %--------------------------------------------------------------------%
%
% Bab II — Kajian Pustaka
%
% @author Petra Novandi and Modified by Kardynan Parulian
%
%--------------------------------------------------------------------%
%
% Berisi kajian literatur, dasar teori, dan studi terkait. Tambahkan
% referensi baru di references.bib dan sitasi dengan \parencite{key}.
%
%--------------------------------------------------------------------%

\chapter{Kajian Pustaka}
\section{Pemrosesan Bahasa Alami}
\textit{Natural Language Processing} (NLP) merupakan cabang kecerdasan buatan yang berfokus pada interaksi antara komputer dan bahasa manusia melalui pemahaman, penafsiran, serta penghasilan teks atau tuturan secara alami \parencite{jurafsky2025speech}. Pendekatan NLP bergeser dari penggunaan aturan linguistik manual menuju metode statistik dan pembelajaran mesin yang memungkinkan komputer mempelajari pola bahasa dari data. Pergeseran ini berlanjut ke \textit{deep learning} atau pembelajaran mendalam, yaitu pembelajaran menggunakan jaringan saraf berlapis yang mempelajari representasi bahasa secara otomatis dari data, sehingga ketergantungan pada rekayasa fitur manual berkurang \parencite{jurafsky2025speech}.

Perkembangan \textit{deep learning} pada pemrosesan bahasa alami (NLP) ditandai dengan hadirnya Transformer, yaitu arsitektur jaringan saraf yang memproses sekuens data secara paralel menggunakan mekanisme \textit{self-attention} untuk mempelajari hubungan kontekstual antarkata \parencite{vaswani2017attention}. Keunggulan arsitektur ini dibandingkan model sebelumnya meliputi kemampuan pemrosesan paralel yang mempersingkat waktu pelatihan, penanganan konteks yang panjang, dan stabilitas pelatihan pada skala besar. Transformer kemudian menjadi dasar bagi pengembangan model bahasa pralatih \parencite{gardazi2025bert}, salah satunya \textit{Bidirectional Encoder Representations from Transformers} (BERT) yang dipublikasikan pada tahun 2018. BERT secara khusus menggunakan bagian \textit{encoder} dan dilatih secara \textit{bidirectional}, sehingga mampu mempertimbangkan konteks teks dari arah kiri maupun kanan secara bersamaan. Proses prapelatihan model ini memanfaatkan dua strategi utama, yakni \textit{Masked Language Model} (MLM) untuk memprediksi token yang ditutupi, dan \textit{Next Sentence Prediction} (NSP) \parencite{devlin2019bert}.

Pendekatan model pralatih tersebut selanjutnya diadaptasi ke dalam bahasa Indonesia melalui IndoBERT, yang dievaluasi kapabilitasnya menggunakan \textit{benchmark} IndoNLU melalui 12 tugas \textit{downstream} \parencite{wilie2020indonlu}. Salah satu tugas \textit{downstream} yang menjadi fokus pada penelitian ini adalah analisis sentimen, yang pada dasarnya merupakan tugas klasifikasi teks dan diselesaikan melalui proses \textit{fine-tuning} \parencite{devlin2019bert, wilie2020indonlu}. Secara konsep, jenis klasifikasi teks dibedakan oleh sifat ruang labelnya yaitu, \textit{binary classification} jika terdapat dua kelas dengan satu label per teks, \textit{multi-class classification} jika terdapat lebih dari dua kelas saling lepas, dan \textit{multi-label classification} jika satu teks dapat menerima lebih dari satu label \parencite{galke2025really}.

Salah satu bentuk dari klasifikasi teks adalah analisis sentimen, yaitu penentuan polaritas teks ke dalam kategori positif, negatif, atau netral \parencite{ferdous2024sentiment}. Kelas netral secara khusus diberikan pada teks yang tidak memuat tendensi positif maupun negatif. Mengingat terdapat tiga kelas yang saling lepas dengan penugasan satu label per teks, tugas analisis sentimen pada penelitian ini dikategorikan sebagai \textit{multi-class classification}.

Lebih spesifik, analisis sentimen dapat diterapkan pada domain teks keuangan yang membahas kondisi atau peristiwa finansial, seperti berita ekonomi maupun unggahan media sosial terkait saham. Pada domain ini, \textcite{malo2014good} merumuskan penilaian teks dari sudut pandang investor—yakni dampak sentimen terhadap harga saham—yang kemudian diadaptasi menjadi skema tiga label pada \textit{Indonesian Financial PhraseBank} \parencite{maharani2023domain}. Teks keuangan memiliki karakteristik unik karena memuat istilah teknis, ekspresi numerik, dan makna sentimen yang sangat bergantung pada konteks \parencite{araci2019finbert, maharani2023domain}. Sejauh ini, pendekatan menggunakan IndoBERT untuk analisis sentimen teks keuangan berbahasa Indonesia telah diterapkan setidaknya dalam empat penelitian terdahulu \parencite{maharani2023domain, jiwandono2026finetuned}.

Meskipun memiliki kinerja yang baik, \textcite{tanjaya2025optimizing} melaporkan adanya indikasi \textit{overfitting} pada penerapan model berbasis Transformer untuk klasifikasi sentimen \textit{tweet} saham di Indonesia. Hal ini dibuktikan dengan [masukkan kalimat bukti: misal, "selisih akurasi sebesar X" antara data latih dan data uji"]. Selain tantangan evaluasi tersebut, ukuran model yang masif juga membatasi penerapannya pada perangkat dengan sumber daya komputasi terbatas \parencite{abdurrahman2023implementing}. Sebagai gambaran, model dasar IndoBERT memiliki [masukkan angka parameter: misal, "sekitar 110 juta"] parameter \parencite{wilie2020indonlu}. Rangkaian kondisi inilah yang memunculkan urgensi penerapan teknik kompresi model, salah satunya melalui kuantisasi, yang akan dieksplorasi lebih lanjut pada subbab berikutnya.

\section{Analisis Sentimen}
Analisis sentimen pada teks keuangan merupakan tugas NLP yang bertujuan mengidentifikasi orientasi afektif dalam teks keuangan, seperti berita pasar, laporan perusahaan, atau komentar investor, ke dalam kategori positif, negatif, atau netral \parencite{maharani2023domain}. Tugas ini bergantung pada domain karena istilah keuangan memiliki makna yang berbeda dari bahasa umum \parencite{araci2019finbert}. Sebagai contoh, \textcite{maharani2023domain} memberi label netral pada kalimat yang hanya menyebutkan durasi dan nilai suatu proyek tanpa penilaian apa pun, walaupun kalimat tersebut memuat angka dan istilah proyek. Selain pendekatan berbasis leksikon \parencite{mishev2020evaluation}, tersedia model pra-latih seperti FinBERT yang dilatih khusus pada korpus keuangan \parencite{yang2020finbert}. Pada bahasa Indonesia, IndoBERT yang di-\textit{fine-tune} pada berita keuangan mencapai akurasi 0,94 dan Macro F1 sebesar 0,71, dengan catatan bahwa label pada data tersebut berupa \textit{pseudo-label} yang dihasilkan mesin, bukan label manusia \parencite{bimasena2026benchmarking}.

Selain label diskret, sentimen dapat diformulasikan sebagai indikator kontinu yang merepresentasikan intensitas sentimen. \textcite{barter2026bondbert} menggunakan \textit{Normalised Difference Index} (NDI) dengan formula $s = (p_+ - p_-)/(p_+ + p_-)$, dengan $s \in [-1, 1]$, untuk mengukur sentimen pada data obligasi pemerintah Inggris. Menurut penulisnya, NDI tidak memberi ruang bagi kelas netral dan menonjolkan polaritas. \textcite{mishev2020evaluation} menggunakan formula serupa, yaitu $\text{Polarity} = (\text{Pos} - \text{Neg})/(\text{Pos} + \text{Neg})$, pada analisis sentimen berbasis leksikon. \textcite{kim2023forecasting} menggunakan pendekatan berbeda, yaitu menjumlahkan label (+1/0/$-$1) yang dikalikan probabilitas label tersebut per hari pada ringkasan artikel \textit{The New York Times}.

Formula yang diadopsi dalam penelitian ini adalah $S = P(\text{Positif}) - P(\text{Negatif})$, dengan $S \in [-1, 1]$. Secara matematis, $S$ merupakan nilai harapan label bernilai +1, 0, dan $-$1 di bawah distribusi probabilitas keluaran model, yaitu $S = \sum_{y \in \{+1,0,-1\}} y \, P(y)$. Bentuk ini sejalan dengan penjumlahan label berbobot probabilitas pada \textcite{kim2023forecasting}, tetapi diterapkan pada tingkat [kalimat/dokumen, sesuaikan] dan bukan agregasi harian. Perbedaan $S$ dengan NDI terletak pada penyebut $p_+ + p_-$.

Ketika penyebut $p_+ + p_-$ bernilai kecil, selisih $p_+ - p_-$ yang kecil diperbesar oleh pembagian. Sebagai ilustrasi, pasangan $(p_+, p_-) = (0; 0{,}001)$ menghasilkan NDI sebesar $-1$, sedangkan $(0{,}001; 0)$ menghasilkan $+1$, padahal selisih probabilitasnya hanya 0,001 dan nilai $S$ masing-masing $-0{,}001$ dan $+0{,}001$. Pada Indonesian Financial PhraseBank yang 61,44\% berlabel netral \parencite{maharani2023domain}, penyebut dapat bernilai kecil pada prediksi netral dengan keyakinan tinggi. Gangguan numerik akibat presisi rendah pada kuantisasi dapat terbaca sebagai pergeseran besar pada NDI yang tidak mencerminkan perubahan perilaku model. Karena itu, penelitian ini memilih $S$ tanpa penyebut, sehingga massa probabilitas netral menekan nilai $S$ ke arah 0 alih-alih dibuang.

Indikator kontinu ini relevan untuk mengevaluasi dampak kuantisasi karena kuantisasi dapat mengubah keluaran model meskipun akurasi nyaris sama. \textcite{sanjaya2026quantization} melaporkan akurasi IndoBERT pada SmSA sebesar [isi FP32] pada FP32 dan [isi] setelah kuantisasi, sedangkan korelasi Spearman untuk stabilitas atribusi turun dari 0,967 pada PTQ-INT8 menjadi sekitar 0,55 pada varian QAT setelah ekspor ONNX. Kesesuaian (\textit{agreement}) prediksi PTQ-INT4 terhadap FP32 sebesar 0,974, yang berarti sekitar 2,6\% prediksi berubah. Temuan ini menunjukkan bahwa perubahan perilaku model akibat kuantisasi tidak terlihat pada akurasi saja, sehingga diperlukan ukuran tambahan yang bekerja pada keluaran probabilitas.

\section{BERT dan IndoBERT}

\subsection{BERT}

BERT diperkenalkan oleh \textcite{devlin2019bert} sebagai model bahasa pra-latih yang menggunakan \textit{encoder} Transformer secara dua arah (\textit{bidirectional}). BERT dilatih dengan dua tugas, yaitu \textit{Masked Language Modeling} (MLM) dan \textit{Next Sentence Prediction} (NSP) \parencite{devlin2019bert}. Pada MLM, sebagian token dalam input disembunyikan dan model diminta memprediksinya berdasarkan konteks kiri dan kanan, sehingga BERT menangkap konteks lebih kaya dibandingkan model bahasa searah \parencite{devlin2019bert}. BERT tersedia dalam beberapa ukuran, seperti BERT-base dengan 110 juta parameter dan BERT-large dengan 340 juta parameter \parencite{devlin2019bert}.

Mekanisme inti Transformer adalah \textit{scaled dot-product attention}, yang menghitung bobot perhatian berdasarkan hubungan antara \textit{query} ($Q$), \textit{key} ($K$), dan \textit{value} ($V$) \parencite{vaswani2017attention}:

\begin{equation}
  \text{Attention}(Q, K, V) = \text{softmax}\left(\frac{QK^T}{\sqrt{d_k}}\right)V
  \label{eq:attention}
\end{equation}

\noindent dengan $d_k$ adalah dimensi \textit{key}. Karena mekanisme \textit{attention} tidak memiliki urutan bawaan, Transformer menambahkan \textit{positional encoding} untuk mempertahankan informasi urutan \parencite{vaswani2017attention}. \textit{Multi-head attention} memperluas \textit{attention} tunggal dengan menjalankan $h$ \textit{attention} secara paralel:

\begin{equation}
  \text{MultiHead}(Q, K, V) = \text{Concat}(\text{head}_1, \dots, \text{head}_h)W^O
  \label{eq:multihead}
\end{equation}

\noindent dengan $\text{head}_i = \text{Attention}(QW_i^Q, KW_i^K, VW_i^V)$, di mana $W_i^Q$, $W_i^K$, dan $W_i^V$ adalah matriks proyeksi untuk \textit{head} ke-$i$ dan $W^O$ adalah matriks proyeksi keluaran \parencite{vaswani2017attention}. Lapisan \textit{feed-forward} dan \textit{layer normalization} melengkapi blok \textit{encoder} \parencite{devlin2019bert}.

Objektif pelatihan MLM dirumuskan sebagai berikut:

\begin{equation}
  \mathcal{L}_{\text{MLM}} = -\sum_{i \in \mathcal{M}} \log P(w_i \mid \mathbf{w}_{\setminus \mathcal{M}})
  \label{eq:mlm}
\end{equation}

\noindent dengan $\mathcal{M}$ adalah himpunan indeks token yang disembunyikan, $w_i$ adalah token target, dan $\mathbf{w}_{\setminus \mathcal{M}}$ adalah konteks token yang tersisa \parencite{devlin2019bert}. Ukuran parameter yang besar membuat BERT sulit dijalankan pada perangkat dengan sumber daya terbatas, sehingga kuantisasi menjadi teknik kompresi yang relevan.

\subsection{IndoBERT}

IndoBERT merupakan adaptasi BERT untuk bahasa Indonesia yang dilatih pada korpus bahasa Indonesia berskala besar \parencite{wilie2020indonlu}. \textcite{wilie2020indonlu} mengembangkannya dengan korpus Indo4B yang terdiri dari sekitar empat miliar kata, mencakup teks formal dan kolokial dari berbagai sumber seperti media sosial, blog, berita, dan situs web. IndoBERT memiliki arsitektur yang mirip dengan BERT, tetapi kosakata dan representasinya disesuaikan dengan karakteristik bahasa Indonesia \parencite{maharani2023domain}. Penelitian ini menggunakan \textit{checkpoint} \textit{indobenchmark/indobert-base-p1}, yaitu varian \textit{base} dengan 124,5 juta parameter, 12 lapisan, 12 \textit{head}, dan \textit{hidden size} 768 \parencite{wilie2020indonlu, maharani2023domain}.

Meskipun performanya baik, ukuran parameter IndoBERT yang besar memerlukan kompresi agar lebih efisien \parencite{abdurrahman2023implementing}. Penelitian terdahulu pada IndoBERT menunjukkan bahwa kuantisasi INT8 pada bobot menghasilkan degradasi performa yang minimal, kuantisasi 4-bit pada bobot masih mempertahankan performa pada tugas klasifikasi tetapi lebih merugikan tugas \textit{sequence labeling}, dan kuantisasi aktivasi 4-bit menyebabkan penurunan yang sangat besar \parencite{abdurrahman2023implementing}. Angka-angka rinci studi tersebut disajikan pada Subbab 2.8.

\section{Kuantisasi Model}

Kuantisasi adalah teknik untuk melakukan komputasi dan menyimpan tensor pada presisi bit yang lebih rendah \parencite{abdurrahman2023implementing}. Secara matematis, kuantisasi merupakan pemetaan nilai riil kontinu ke himpunan nilai diskret berhingga, dengan input di luar rentang yang dipilih dipotong (\textit{clipped}) agar berada dalam batas rentang tersebut \parencite{gholami2021survey}. Model terkuantisasi dapat menjalankan sebagian atau seluruh operasinya menggunakan presisi rendah \parencite{wu2023understanding}. Dengan pendekatan ini, kebutuhan memori dan komputasi dapat dikurangi, tetapi besarnya degradasi performa bergantung pada presisi dan cakupan kuantisasi \parencite{gholami2021survey}.

Presisi yang umum digunakan dalam kuantisasi adalah FP32, FP16, INT8, dan INT4. Notasi W$x$A$y$ menunjukkan bobot $x$-bit dan aktivasi $y$-bit; W8A8 berarti bobot dan aktivasi 8-bit, sedangkan W4A4 berarti keduanya 4-bit \parencite{wu2023understanding}. Dalam penelitian ini, W4A16 menunjukkan bobot 4-bit dengan aktivasi 16-bit yang tidak dikuantisasi (skenario \textit{weight-only}), dan W4A8 menunjukkan bobot 4-bit dengan aktivasi 8-bit. Kuantisasi dapat diterapkan pada bobot saja (\textit{weight-only}) atau pada bobot dan aktivasi sekaligus (\textit{weight-activation}) \parencite{wu2023understanding}. Kuantisasi bobot dan aktivasi memberikan penghematan lebih besar, tetapi memerlukan kalibrasi yang lebih cermat untuk menghindari degradasi performa \parencite{gholami2021survey}. Relatif terhadap model FP32, model berpresisi 16-bit berukuran sekitar 50\% dan model 8-bit sekitar 25\% \parencite{abdurrahman2023implementing}.

Kuantisasi simetris untuk bilangan bulat $b$-bit dengan pendekatan pembagi didefinisikan sebagai berikut:

\begin{equation}
  x_q = \text{clip}\left(\text{round}\left(\frac{x}{s}\right), -2^{b-1}, 2^{b-1}-1\right)
  \label{eq:kuantisasi}
\end{equation}

\noindent dengan faktor skala:

\begin{equation}
  s = \frac{\alpha}{2^{b-1}-1}
  \label{eq:skala}
\end{equation}

\noindent Pada Persamaan~\eqref{eq:kuantisasi} dan \eqref{eq:skala}, $x$ adalah nilai riil input, $x_q$ adalah nilai terkuantisasi, $b$ adalah jumlah bit representasi, dan $\alpha$ adalah nilai absolut maksimum dari rentang yang dipilih untuk kuantisasi \parencite{gholami2021survey}. Karena skema ini simetris terhadap nol, \textit{zero point} bernilai nol. Operasi \textit{clip} memastikan nilai berada dalam rentang bilangan bulat bertanda. Dekuantisasi dilakukan dengan mengalikan kembali nilai terkuantisasi dengan faktor skala, yaitu $\hat{x} = x_q \cdot s$ \parencite{gholami2021survey}.

\section{Post-Training Quantization}

\textit{Post-Training Quantization} (PTQ) adalah proses di mana skala dan \textit{zero point} untuk setiap tensor dalam model \textit{neural network} ditentukan terlebih dahulu dengan menjalankan model untuk mendapatkan statistik setiap lapisan \parencite{abdurrahman2023implementing}. PTQ tidak memerlukan \textit{retraining}, sehingga lebih praktis untuk \textit{deployment} model yang sudah dilatih \parencite{gholami2021survey}. Data kalibrasi diperlukan ketika aktivasi dikuantisasi, karena data tersebut digunakan untuk mengestimasi statistik aktivasi per lapisan \parencite{williams2024impact}. \textcite{williams2024impact} menunjukkan bahwa pemilihan data kalibrasi memengaruhi kualitas kuantisasi pada model bahasa, sedangkan \textcite{hubara2021accurate} menemukan bahwa PTQ dengan set kalibrasi kecil masih dapat mencapai akurasi kompetitif pada model \textit{vision}.

Penelitian ini menggunakan skema kuantisasi simetris dengan bobot dikuantisasi per-\textit{channel} berbasis \textit{absmax}. Faktor skala per-\textit{channel} dirumuskan sebagai berikut:

\begin{equation}
  s_j = \frac{\max(|W_j|)}{2^{b-1}-1}
  \label{eq:perchannel}
\end{equation}

\noindent dengan $W_j$ adalah bobot pada \textit{channel} ke-$j$ dan $s_j$ adalah faktor skala untuk \textit{channel} tersebut \parencite{gholami2021survey}. Bobot hasil rekonstruksi diperoleh dengan:

\begin{equation}
  \hat{W}_j = s_j \cdot \text{round}\left(\frac{W_j}{s_j}\right)
  \label{eq:dekuantisasi}
\end{equation}

\noindent dengan $\hat{W}_j$ adalah bobot hasil dekuantisasi \parencite{gholami2021survey}. Persamaan~\eqref{eq:dekuantisasi} menunjukkan bahwa bobot asli didekati kembali dari representasi bilangan bulat dengan mengalikan faktor skala. Perlakuan kuantisasi pada setiap kondisi percobaan, termasuk kondisi yang tidak mengkuantisasi aktivasi, dijelaskan pada Bab III.

Degradasi performa pada PTQ bergantung pada cakupan kuantisasi. \textcite{abdurrahman2023implementing} melaporkan bahwa kuantisasi aktivasi 4-bit jauh lebih merusak performa daripada kuantisasi bobot 4-bit pada IndoBERT, yang menunjukkan bahwa aktivasi lebih sensitif terhadap kuantisasi agresif dibandingkan bobot. \textcite{sanjaya2026quantization} juga melaporkan bahwa PTQ \textit{weight-only} INT4 hampir tidak menurunkan metrik klasifikasi, tetapi memengaruhi stabilitas atribusi dan menyebabkan sebagian prediksi berubah (lihat Subbab 2.2).

Berbeda dengan PTQ, \textit{Quantization-Aware Training} (QAT) memasukkan operasi kuantisasi ke dalam model sebelum pelatihan atau \textit{fine-tuning} \parencite{wu2023understanding}. QAT memungkinkan model beradaptasi dengan bobot dan aktivasi terkuantisasi sehingga degradasi performa dapat ditekan \parencite{wu2023understanding}. Q8BERT melaporkan bahwa kuantisasi 8-bit dengan QAT menghasilkan degradasi akurasi sekitar 1\% pada beberapa tugas GLUE \parencite{zafrir2019q8bert}. \textcite{wu2023understanding} melaporkan bahwa QAT dengan \textit{knowledge distillation} pada BERT-base tidak menyebabkan degradasi akurasi pada tugas klasifikasi GLUE, sedangkan pada model \textit{decoder-only} seperti GPT2 terjadi penurunan \textit{perplexity} yang signifikan. Namun, QAT memerlukan biaya komputasi lebih tinggi karena memerlukan \textit{retraining} tambahan. \textcite{sanjaya2026quantization} menemukan bahwa QAT mengubah atribusi model bahkan pada presisi penuh, dengan korelasi Spearman turun menjadi 0,788 sebelum kuantisasi diterapkan, yang menunjukkan bahwa \textit{fine-tuning} tambahan itu sendiri mengubah perilaku internal model. Penelitian ini memilih PTQ untuk mengevaluasi kelayakan kompresi tanpa \textit{retraining}.

\section{Dataset Finansial Bahasa Indonesia}

\subsection{IndoFinSent}

IndoFinSent merupakan dataset analisis sentimen finansial bahasa Indonesia asli yang dibuat dengan \textit{scraping} judul artikel berita keuangan dari berbagai sumber seperti Liputan6, CNBC, Detik, Sindo, Jawapos, dan Metrotv News \parencite{maharani2023domain}. Setiap kalimat dilabeli secara manual oleh dua anotator dan diperiksa silang \parencite{maharani2023domain}. Dataset ini berisi 2.274 entri dengan distribusi label positif 52,77\%, netral 26,25\%, dan negatif 20,98\% \parencite{maharani2023domain}. Split resmi dataset terdiri dari 1.591 data latih, 227 data validasi, dan 456 data uji \parencite{maharani2023domain}. Karena dibuat dari teks asli bahasa Indonesia, dataset ini mempertahankan nuansa bahasa, referensi budaya, dan ekspresi idiomatik yang spesifik \parencite{maharani2023domain}, dan pelabelan manual oleh dua anotator umumnya memberi kualitas label yang lebih baik daripada pelabelan otomatis.

Satu keterbatasan perlu dicatat, yaitu IndoFinSent berisi judul berita yang merujuk pada satu institusi keuangan tertentu (contohnya BRI), sehingga generalisasi hasil ke institusi lain atau domain finansial yang lebih luas harus dilakukan dengan hati-hati \parencite{maharani2023domain}. Selain itu, distribusi kelas yang didominasi positif membuat kelas minoritas (negatif 20,98\%) berisiko kurang terwakili dalam pelatihan. Karena itu, Macro F1-Score menjadi metrik utama karena memberikan bobot setara pada setiap kelas \parencite{hinojosalee2024performance}, dilengkapi evaluasi per kelas untuk melihat apakah kuantisasi lebih merugikan kelas minoritas.

\subsection{Indonesian Financial PhraseBank}

Indonesian Financial PhraseBank merupakan versi terjemahan Financial PhraseBank ke dalam bahasa Indonesia menggunakan Google Translate \parencite{maharani2023domain}. Dataset ini memakai subset ``\textit{all agree}'' yang terdiri dari 2.264 kalimat keuangan beserta label sentimennya \parencite{malo2014good}. Sebagian besar kalimat berlabel netral (61,44\%), diikuti positif (25,18\%) dan negatif (13,38\%) \parencite{maharani2023domain}. Split resmi dataset terdiri dari 1.629 data latih, 182 data validasi, dan 453 data uji \parencite{maharani2023domain}.

Penerjemahan memungkinkan perbandingan dengan penelitian terdahulu yang menggunakan Financial PhraseBank versi Inggris \parencite{araci2019finbert}. Terjemahan berpotensi menimbulkan pergeseran makna, tetapi subset ``\textit{all agree}'' memastikan hanya kalimat dengan kesepakatan anotator yang tinggi yang disertakan \parencite{malo2014good}. Dominasi kelas netral membawa konsekuensi bagi pengukuran \textit{score drift}: sebagian besar nilai $S = P(\text{Positif}) - P(\text{Negatif})$ diperkirakan mendekati 0, sehingga MAE \textit{drift} secara keseluruhan berpotensi tampak kecil. Karena itu, \textit{score drift} perlu dilaporkan per kelas prediksi dan dilengkapi KL \textit{divergence} pada distribusi tiga kelas (Subbab 2.7.2).

\section{Metrik Evaluasi}

\subsection{Klasifikasi Diskret}

\textit{Macro F1-Score} dihitung dengan menghitung F1-score untuk setiap kelas secara terpisah, kemudian mengambil rata-ratanya \parencite{hinojosalee2024performance}. Metrik ini memberikan bobot setara pada setiap kelas sehingga lebih peka terhadap ketidakseimbangan kelas dibandingkan \textit{Micro F1-Score} \parencite{hinojosalee2024performance}. Rumusnya adalah sebagai berikut:

\begin{equation}
  F1_{\text{macro}} = \frac{1}{N} \sum_{i=1}^{N} F1_i
  \label{eq:macrof1}
\end{equation}

\noindent dengan $N$ adalah jumlah kelas dan $F1_i$ adalah F1-score untuk kelas ke-$i$ \parencite{hinojosalee2024performance}.

Akurasi untuk klasifikasi tiga kelas dihitung sebagai proporsi prediksi yang benar terhadap seluruh data uji:

\begin{equation}
  \text{Akurasi} = \frac{\sum_{i=1}^{N} \text{TP}_i}{N_{\text{total}}}
  \label{eq:akurasi}
\end{equation}

\noindent dengan $\text{TP}_i$ adalah jumlah prediksi benar untuk kelas ke-$i$ (elemen diagonal matriks konfusi) dan $N_{\text{total}}$ adalah jumlah total sampel uji. Akurasi mudah dipahami tetapi dapat menyesatkan pada data yang tidak seimbang antar-kelas \parencite{sokolova2009systematic}. Contohnya adalah studi \textcite{bimasena2026benchmarking} yang melaporkan akurasi 0,94 bersama Macro F1 hanya 0,71 pada data berita keuangan. \textit{Precision} dan \textit{recall} juga dilaporkan per kelas untuk melihat apakah kuantisasi lebih merugikan kelas minoritas.

\subsection{Definisi Score Drift}

\textit{Score drift} didefinisikan sebagai selisih antara nilai indikator sentimen kontinu $S$ pada model terkuantisasi dan model FP32. Secara per sampel, \textit{drift} dirumuskan sebagai:

\begin{equation}
  d_i = S_{q,i} - S_{\text{FP32},i}
  \label{eq:drift}
\end{equation}

\noindent dengan $q$ menunjukkan presisi terkuantisasi dan tanda negatif menunjukkan bahwa model terkuantisasi lebih pesimis. Urutan pengurangan tidak memengaruhi MAE tetapi memengaruhi tanda bias, sehingga penelitian ini konsisten menggunakan $S_{q} - S_{\text{FP32}}$. \textit{Baseline} adalah model FP32 hasil \textit{fine-tuning} yang sama, pada \textit{test set} dan tokenisasi yang sama.

Secara agregat, \textit{score drift} diukur menggunakan \textit{Mean Absolute Error} (MAE), korelasi Pearson, dan korelasi Spearman. MAE dirumuskan sebagai:

\begin{equation}
  \text{MAE} = \frac{1}{n} \sum_{i=1}^{n} |d_i|
  \label{eq:mae}
\end{equation}

Korelasi Pearson antara nilai $S$ sebelum dan sesudah kuantisasi dirumuskan sebagai:

\begin{equation}
  r = \frac{\sum_{i=1}^{n} (S_{\text{FP32},i} - \bar{S}_{\text{FP32}})(S_{q,i} - \bar{S}_{q})}{\sqrt{\sum_{i=1}^{n} (S_{\text{FP32},i} - \bar{S}_{\text{FP32}})^2} \sqrt{\sum_{i=1}^{n} (S_{q,i} - \bar{S}_{q})^2}}
  \label{eq:pearson}
\end{equation}

Korelasi Spearman dihitung dengan menerapkan korelasi Pearson pada peringkat nilai $S$ alih-alih nilai aslinya, sehingga mengukur apakah urutan peringkat sampel dipertahankan setelah kuantisasi.

\textit{Label flip rate} didefinisikan sebagai proporsi sampel yang label prediksinya berubah antara model FP32 dan model terkuantisasi:

\begin{equation}
  \text{Flip Rate} = \frac{1}{n} \sum_{i=1}^{n} \mathbb{1}[\hat{y}_{\text{FP32},i} \neq \hat{y}_{q,i}]
  \label{eq:fliprate}
\end{equation}

\noindent dengan $\hat{y}_{\text{FP32},i}$ adalah label prediksi model FP32 untuk sampel ke-$i$, $\hat{y}_{q,i}$ adalah label prediksi model terkuantisasi, dan $\mathbb{1}[\cdot]$ adalah fungsi indikator.

Selain itu, KL \textit{divergence} pada distribusi tiga kelas (positif, netral, negatif) digunakan untuk mengukur seberapa jauh distribusi probabilitas keluaran bergeser:

\begin{equation}
  D_{\text{KL}}(P_{\text{FP32}} \,\|\, P_q) = \sum_{c=1}^{3} \tilde{P}_{\text{FP32}}(c) \ln \frac{\tilde{P}_{\text{FP32}}(c)}{\tilde{P}_q(c)}
  \label{eq:kl}
\end{equation}

\noindent dengan $\tilde{P}(c) = (P(c) + \varepsilon)/(1 + 3\varepsilon)$ adalah probabilitas yang dihaluskan dengan $\varepsilon = 10^{-8}$ agar tidak terjadi pembagian dengan nol, dan logaritma yang digunakan adalah logaritma natural. Prosedur pengukuran, termasuk pengujian statistik dan analisis hubungan \textit{drift} dengan margin keputusan model, dijelaskan pada Bab III.

\subsection{Efisiensi Komputasi}

Ukuran model diukur dalam megabyte (MB) dan dibandingkan antara presisi FP32, INT8, dan INT4. Estimasi teoretis dari jumlah bit dapat berbeda dari ukuran aktual berkas. \textcite{sanjaya2026quantization} melaporkan bahwa PTQ-INT8 mengurangi ukuran model IndoBERT dari 475 MB menjadi 230 MB, atau sekitar 52\%, sedangkan PTQ-INT4 \textit{weight-only} tetap berukuran 230 MB karena disimpan sebagai INT8 dengan metadata tambahan.

Latensi diukur dalam milidetik per sampel (ms/sampel) dan \textit{throughput} dalam sampel per detik. \textcite{sanjaya2026quantization} melaporkan bahwa PTQ-INT8 mengurangi latensi dari 25,1 ms menjadi 9,4 ms pada CPU, atau sekitar 2,7$\times$ lebih cepat, sedangkan PTQ-INT4 \textit{weight-only} justru lebih lambat (32,7 ms) karena dekuantisasi ke FP32 sebelum komputasi. \textcite{wu2023understanding} melaporkan bahwa \textit{pipeline} INT4 dengan QAT dan \textit{knowledge distillation} mencapai 8,5$\times$ lebih cepat untuk skenario berorientasi latensi dan 3$\times$ lebih cepat untuk \textit{throughput} dibandingkan FP16 pada GPU A6000. Namun, percepatan INT4 memerlukan dukungan \textit{kernel hardware} tertentu dan tidak terwujud pada \textit{weight-only quantization} tanpa dukungan tersebut \parencite{wu2023understanding}.

Karena simulasi kuantisasi (\textit{fake quantization}) tidak menghasilkan percepatan, penelitian ini melaporkan ukuran model dan estimasi \textit{bit-operations} (BOPs) sebagai proksi efisiensi utama, serta latensi sebagai metrik pelengkap dengan catatan bahwa hasilnya bergantung pada dukungan \textit{hardware} dan \textit{kernel}. Ukuran model dilaporkan dalam dua bentuk, yaitu ukuran teoretis berdasarkan jumlah bit dan ukuran aktual berkas, karena penyimpanan INT4 dapat dikemas sebagai INT8 \parencite{sanjaya2026quantization}.

\section{Penelitian Relevan}

Penelitian terdahulu yang relevan dan terbit dalam lima tahun terakhir dirangkum pada Tabel~\ref{tab:penelitian-terdahulu}. Studi-studi tersebut dapat dikelompokkan ke dalam tiga jalur yang saling melengkapi, yaitu sentimen finansial berbahasa Indonesia, kuantisasi model bahasa, dan indikator sentimen kontinu.

Pada jalur sentimen finansial Indonesia, \textcite{maharani2023domain} melakukan \textit{post-training} IndoBERT pada korpus finansial Indonesia sebesar 647.586 kata. Pada Financial PhraseBank terjemahan, model \textit{base} mencapai F1 0,91 pada \textit{baseline} dan 0,94 setelah \textit{post-training} dengan 100\% data latih. Pada IndoFinSent, model \textit{base} yang di-\textit{post-training} mencapai F1 0,79--0,81, sedangkan model \textit{large} mencapai F1 0,78--0,79 \parencite{maharani2023domain}. \textcite{tanjaya2025optimizing} melakukan \textit{fine-tuning} IndoBERTweet pada 3.288 tweet saham berbahasa Indonesia dan melaporkan F1 makro 0,8177 pada data uji serta 0,8106 pada data uji dengan \textit{noise} sintetis, dengan indikasi \textit{overfitting} yang diakui penulisnya. \textcite{bimasena2026benchmarking} membandingkan IndoBERT dan Multilingual BERT untuk klasifikasi sentimen berita keuangan Indonesia, dengan IndoBERT mencapai akurasi 0,94 dan Macro F1 0,71. Data \textcite{bimasena2026benchmarking} menggunakan \textit{pseudo-label} yang dihasilkan mesin, bukan \textit{ground truth} manusia. Ketiga studi ini melaporkan performa IndoBERT untuk sentimen finansial Indonesia, tetapi semuanya memakai metrik klasifikasi diskret dan tidak mengevaluasi dampak kompresi model terhadap keluaran probabilitas.

Pada jalur kuantisasi, \textcite{sanjaya2026quantization} mengevaluasi PTQ dan QAT pada IndoBERT untuk klasifikasi sentimen tiga kelas SmSA pada presisi FP16, INT8, dan INT4. Akurasi hampir tidak berubah (92,9--93,2\%), sedangkan stabilitas atribusi bergantung pada jalur kompresi: PTQ-INT8 mempertahankan korelasi Spearman 0,967 terhadap FP32, sementara varian QAT turun hingga sekitar 0,55 setelah ekspor ONNX. Studi tersebut memakai data umum (SmSA), INT4 hanya pada bobot, dan metrik berbasis atribusi. Studi tersebut tidak mendefinisikan indikator sentimen kontinu dari probabilitas kelas maupun mengukur pergeserannya per sampel akibat kuantisasi, dan tidak dilakukan pada domain finansial. \textcite{wu2023understanding} meneliti kuantisasi INT4 pada bobot dan aktivasi (W4A4) pada model Transformer dengan QAT dan \textit{knowledge distillation}. Mereka menemukan bahwa W4A4 tidak menyebabkan degradasi akurasi pada model \textit{encoder-only} seperti BERT, tetapi menyebabkan penurunan besar pada model \textit{decoder-only}. \textit{Pipeline} INT4 mereka dilaporkan 8,5$\times$ lebih cepat untuk skenario berorientasi latensi dan 3$\times$ lebih cepat untuk \textit{throughput} dibandingkan FP16. \textcite{williams2024impact} menunjukkan bahwa pemilihan data kalibrasi memengaruhi performa \textit{downstream} model terkuantisasi, dengan tingkat sensitivitas yang berbeda antar metode PTQ dan \textit{pruning}. \textcite{abdurrahman2023implementing} meneliti kuantisasi pada IndoBERT untuk tujuh tugas \textit{downstream} IndoNLU. PTQ INT8 pada bobot mencapai degradasi minimal (rata-rata F1 77,04 terhadap 77,35). PTQ 4-bit pada bobot masih mempertahankan performa pada tugas klasifikasi dengan penurunan F1 sekitar 1 poin, tetapi tugas \textit{sequence labeling} mengalami degradasi lebih besar (FACQA 54,16 menjadi 43,90; NERGRIT 70,92 menjadi 60,31). PTQ 4-bit pada aktivasi menyebabkan penurunan besar, dengan rata-rata F1 turun dari 76,98 pada W16A32 menjadi 17,71 pada W16A4, sedangkan QAT pada aktivasi 4-bit masih menghasilkan rata-rata F1 sebesar 47,29. Pada tugas sentimen SmSA, PTQ W4A32 mencapai F1 85,20, lebih tinggi daripada QAT W4A32 yang mencapai 79,21. Hasil ini menjadi salah satu alasan penelitian ini memilih PTQ, meskipun berasal dari satu dataset dan satu konfigurasi sehingga perlu diuji ulang pada domain finansial.

Pada jalur indikator sentimen kontinu, \textcite{barter2026bondbert} menggunakan NDI pada data obligasi pemerintah Inggris untuk mengukur sentimen yang selaras dengan pergerakan harga obligasi. \textcite{kim2023forecasting} menggunakan skor sentimen kontinu dari probabilitas label \textit{argmax} untuk memprediksi indeks S\&P 500 dan melaporkan bahwa penambahan skor sentimen menurunkan MSE dari 0,0026 menjadi 0,0016 serta menaikkan R$^2$ dari 0,8339 menjadi 0,8950. \textcite{mishev2020evaluation} menggunakan formula polaritas berbasis leksikon pada berita keuangan berbahasa Inggris. Ketiga studi ini menunjukkan kegunaan indikator sentimen kontinu, tetapi tidak ada yang mengukur pergeseran nilai tersebut akibat kuantisasi, dan tidak ada yang mencakup bahasa Indonesia.

{\setlength{\tabcolsep}{3pt}
\captionsetup[longtable]{font=normalsize}
\begingroup\scriptsize
\begin{longtable}{|c|>{\raggedright\arraybackslash}p{1.5cm}|>{\raggedright\arraybackslash}p{1.4cm}|>{\raggedright\arraybackslash}p{2.1cm}|>{\raggedright\arraybackslash}p{1.9cm}|>{\raggedright\arraybackslash}p{1.5cm}|>{\raggedright\arraybackslash}p{3.0cm}|}
  \caption{Penelitian Terdahulu}
  \label{tab:penelitian-terdahulu} \\
  \hline
  \textbf{No} & \textbf{Referensi} & \textbf{Model} & \textbf{Dataset \& Domain} & \textbf{Teknik} & \textbf{Metrik} & \textbf{Hasil Utama} \\ \hline
  \endfirsthead
  \multicolumn{7}{c}{\tablename~\thetable{} (lanjutan)} \\ \hline
  \textbf{No} & \textbf{Referensi} & \textbf{Model} & \textbf{Dataset \& Domain} & \textbf{Teknik} & \textbf{Metrik} & \textbf{Hasil Utama} \\ \hline
  \endhead
  \hline
  \endlastfoot
  1 & Maharani dkk. (2023) & IndoBERT-base (124,5M), IndoBERT-large (335,2M) & Korpus finansial Indonesia (647.586 kata), IndoFinSent (2.274 entri), Financial PhraseBank ID (2.264 entri) & \textit{Post-training} MLM pada korpus finansial Indonesia & F1 & PhraseBank ID: F1 0,91 $\rightarrow$ 0,94 (\textit{post-trained}, 100\% data); IndoFinSent: F1 0,79--0,81 (\textit{base}), 0,78--0,79 (\textit{large}) \\ \hline
  2 & Abdur\-rahman dkk. (2023) & IndoBERT & IndoNLU (7 tugas), termasuk SmSA & PTQ dan QAT; 4-bit dan 8-bit pada bobot dan aktivasi & F1, FLOPs, estimasi ukuran model & PTQ INT8 pada bobot degradasi minimal (77,04 vs 77,35); PTQ INT4 pada bobot turun $\sim$1 poin pada klasifikasi, tetapi \textit{sequence labeling} turun besar (FACQA 54,16 $\rightarrow$ 43,90); PTQ aktivasi 4-bit \textit{collapse} (17,71); QAT pada bobot efektif; QAT aktivasi 4-bit masih 47,29 \\ \hline
  3 & Wu dkk. (2023) & BERT-base (\textit{encoder}), GPT2 (\textit{decoder}) & Inggris, GLUE, Wikitext & QAT + KD; INT4 W4A4 & Akurasi, \textit{perplexity}, latensi, \textit{throughput} & QAT W4A4 tidak mendegradasi \textit{encoder-only}; \textit{decoder-only} turun signifikan; percepatan 8,5$\times$ (latensi) dan 3$\times$ (\textit{throughput}) terhadap FP16 pada GPU A6000 \\ \hline
  4 & Williams \& Aletras (2024) & LLaMA, Vicuna, OPT (7B--33B) & Lima sumber data kalibrasi (C4, CNN-DM, RedPajama, RefinedWeb, Wikipedia) & PTQ (GPTQ, SpQR) dan \textit{pruning} (SparseGPT, Wanda) & Akurasi \textit{zero-shot}, \textit{perplexity} & Pemilihan data kalibrasi memengaruhi performa \textit{downstream}; SpQR paling tidak sensitif, SparseGPT paling sensitif \\ \hline
  5 & Tanjaya dkk. (2025) & Indo\-BERT\-weet & 3.288 tweet saham Indonesia (ID-SMSA) & \textit{Fine-tuning} dengan variasi \textit{preprocessing} dan \textit{hyperparameter} & Macro F1 & F1 makro 0,8177 pada data uji; 0,8106 pada data uji ber-\textit{noise}; indikasi \textit{overfitting} \\ \hline
  6 & Sanjaya dkk. (2026) & IndoBERT-base-p2, XLM-R, mBERT & Indonesia, SmSA, NusaX-senti & PTQ dan QAT; FP16/\allowbreak INT8/\allowbreak INT4 (INT4 hanya bobot) & Akurasi, F1, ECE, \textit{agreement}, korelasi atribusi & Akurasi stabil (92,9--93,2\%); PTQ-INT8 $\rho$ = 0,967; PTQ-INT4 $\rho$ = 0,850; QAT turun hingga $\rho \approx$ 0,55; PTQ-INT8 230 MB dan 9,4 ms; PTQ-INT4 32,7 ms \\ \hline
  7 & Barter dkk. (2026) & BondBERT (FinBERT di-\textit{fine-tune} pada korpus obligasi) & 30.000 artikel obligasi pemerintah Inggris (2018--2025) & \textit{Fine-tuning} dengan label GPT-4.1-nano; NDI & Korelasi dengan \textit{return} obligasi, akurasi arah, \textit{normalised} RMSE, IC, uji Diebold--Mariano & BondBERT konsisten menghasilkan korelasi positif dengan \textit{return} obligasi dan akurasi arah lebih tinggi daripada FinBERT, FinGPT, dan Instruct-FinGPT \\ \hline
  8 & Bimasena dkk. (2026) & IndoBERT, mBERT & 3.366 artikel berita finansial Indonesia (\textit{pseudo-label}) & \textit{Fine-tuning} dengan \textit{class weighting} dan \textit{stratified splitting} & Akurasi, Macro F1 & IndoBERT: akurasi 0,94, Macro F1 0,71; mBERT: akurasi 0,93, Macro F1 0,70 \\ \hline
  9 & \textbf{Penelitian ini} & \textbf{IndoBERT-base-p1} & \textbf{IndoFinSent (2.274 entri), Financial PhraseBank ID (2.264 entri)} & \textbf{PTQ: FP32, INT8-W8A8, INT4-W4A16, INT4-W4A8, INT4-W4A4} & \textbf{Macro F1, akurasi, \textit{score drift} (MAE, Pearson, Spearman, \textit{label flip rate}, KL \textit{divergence}), ukuran model, BOPs, latensi} & \textbf{Dilaporkan pada Bab IV} \\ \hline
\end{longtable}
\endgroup
}

Berdasarkan Tabel~\ref{tab:penelitian-terdahulu}, dapat diidentifikasi empat celah utama. Pertama, studi pada jalur sentimen finansial Indonesia \parencite{maharani2023domain, tanjaya2025optimizing, bimasena2026benchmarking} belum mengevaluasi dampak kompresi model terhadap keluaran probabilitas. Kedua, studi pada jalur kuantisasi \parencite{abdurrahman2023implementing, wu2023understanding, williams2024impact, sanjaya2026quantization} belum mencakup domain finansial Indonesia, dan hanya \textcite{sanjaya2026quantization} yang mengevaluasi kuantisasi pada IndoBERT untuk tugas sentimen meskipun pada data umum (SmSA) dan dengan kuantisasi INT4 hanya pada bobot. Ketiga, studi pada jalur indikator sentimen kontinu \parencite{mishev2020evaluation, kim2023forecasting, barter2026bondbert} belum mengukur pergeseran indikator tersebut akibat kuantisasi. Keempat, tidak ada studi yang secara bersamaan mencakup keempat dimensi: PTQ, sentimen, domain finansial, dan bahasa Indonesia. Penelitian ini mengisi celah tersebut dengan mengevaluasi PTQ pada IndoBERT menggunakan dua dataset finansial Indonesia, yaitu IndoFinSent dan Indonesian Financial PhraseBank, pada lima kondisi: FP32, INT8-W8A8, INT4-W4A16, INT4-W4A8, dan INT4-W4A4.

    %--------------------------------------------------------------------%
%
% Bab III — Analisis dan Perancangan
%
% @author Petra Novandi
%
%--------------------------------------------------------------------%
%
% Berisi analisis masalah, solusi umum, dan rancangan solusi.
% Ganti teks placeholder dengan isi sebenarnya.
%
%--------------------------------------------------------------------%

\chapter{Analisis dan Perancangan}

Bab ini menjelaskan bagaimana persoalan yang dirumuskan di Bab I
dianalisis dan dirancang solusinya. Pembaca seharusnya memahami
pendekatan yang diambil serta alasan pemilihan solusi sebelum
membaca bab evaluasi.

\section{Analisis Masalah}

Jelaskan persoalan tugas akhir secara rinci berdasarkan rumusan
masalah. Uraikan kondisi yang ada, kebutuhan pengguna atau sistem,
serta kendala yang menjadi fokus penyelesaian.

\blindtext%

\section{Solusi Umum}

Deskripsikan solusi atau pendekatan umum yang relevan dengan
persoalan. Bandingkan alternatif yang ada dan jelaskan mengapa
pendekatan tertentu dipilih untuk dikembangkan lebih lanjut.

\blindtext%

\section{Rancangan Solusi}

Uraikan rancangan solusi yang diusulkan, meliputi arsitektur,
komponen, alur kerja, model data, antarmuka, atau aspek teknis
lainnya sesuai topik tugas akhir. Contoh berikut menampilkan data
dari berkas CSV di direktori \texttt{resources/}.

\begin{table}[ht]
    \centering
    \caption{Kebutuhan fungsional sistem}
    \label{tab:sample-requirements}
    \csvautotabular{resources/chapter-3-sample-requirements.csv}
\end{table}

\begin{table}[ht]
    \centering
    \caption{Komponen arsitektur sistem}
    \label{tab:sample-components}
    \csvautotabular{resources/chapter-3-sample-components.csv}
\end{table}

    \subsection{Subbab}

    \blindtext%

    \subsubsection{Subsubbab}

    \blindtext%

\blindtext%

    %--------------------------------------------------------------------%
%
% Bab IV — Evaluasi dan Pembahasan
%
% @author Petra Novandi
%
%--------------------------------------------------------------------%
%
% Berisi tujuan pengujian, skenario, hasil pengujian, dan pembahasan.
% Ganti teks placeholder dengan isi sebenarnya.
%
%--------------------------------------------------------------------%

\chapter{Evaluasi dan Pembahasan}

Bab ini mempresentasikan pengujian solusi yang dirancang pada Bab
III dan membahas apakah tujuan serta rumusan masalah telah
terpenuhi. Hasil yang ditampilkan harus dapat dipertanggungjawabkan
secara objektif.

\section{Tujuan Pengujian}

Tuliskan tujuan pengujian yang spesifik dan terukur. Setiap tujuan
sebaiknya berkaitan langsung dengan tujuan tugas akhir dan
rumusan masalah di Bab I.

\blindtext%

\section{Skenario Pengujian}

Jelaskan metode, lingkungan, data uji, dan langkah-langkah
pengujian yang dilakukan. Skenario harus cukup jelas agar pengujian
dapat direplikasi oleh pembaca.

\blindtext%

\section{Hasil Pengujian}

Sajikan hasil pengujian secara faktual menggunakan tabel, grafik,
atau angka. Hindari interpretasi panjang di subbab ini; fokuskan
pada data dan temuan yang diperoleh.

    \subsection{Subbab}

    \blindtext%

    \subsubsection{Subsubbab}

    \blindtext%

\blindtext%

\section{Pembahasan}

Interpretasikan hasil pengujian dan kaitkan dengan rumusan masalah
serta tujuan tugas akhir. Jelaskan apakah solusi yang diusulkan
berhasil, keterbatasannya, dan perbandingannya dengan studi terkait
jika relevan.

\blindtext%

    %--------------------------------------------------------------------%
%
% Bab V — Penutup
%
% @author Petra Novandi
%
%--------------------------------------------------------------------%
%
% Berisi kesimpulan dan saran. Ganti teks placeholder dengan isi
% sebenarnya.
%
%--------------------------------------------------------------------%

\chapter{Penutup}

Bab ini merangkum hasil tugas akhir dan memberikan saran untuk
pengembangan lebih lanjut. Kesimpulan harus menjawab rumusan masalah
dan tujuan yang telah dirumuskan di awal laporan.

\section{Kesimpulan}

Tuliskan kesimpulan utama berdasarkan hasil analisis, perancangan,
dan evaluasi. Setiap poin kesimpulan sebaiknya merujuk pada tujuan
atau rumusan masalah yang telah diselesaikan.

\blindtext%

\section{Saran}

Berikan saran yang konstruktif untuk pengembangan solusi, perbaikan
metodologi, atau penelitian lanjutan. Saran dapat mencakup aspek
teknis, evaluasi, maupun implementasi di lingkungan nyata.

\blindtext%


    %------------------------------------------------------------%
    % Lampiran & riwayat hidup
    %------------------------------------------------------------%
    \appendix

    %--------------------------------------------------------------------%
%
% Lampiran A — Instrumen Pengujian
%
% @author Petra Novandi
%
%--------------------------------------------------------------------%
%
% Berisi instrumen atau alat bantu pengujian. Ganti teks placeholder
% dengan isi sebenarnya.
%
%--------------------------------------------------------------------%

\chapter{Instrumen Pengujian}

\blindtext%

    %--------------------------------------------------------------------%
%
% Lampiran B — Rincian Kasus Uji
%
% @author Petra Novandi
%
%--------------------------------------------------------------------%
%
% Berisi rincian kasus uji. Ganti teks placeholder dengan isi
% sebenarnya.
%
%--------------------------------------------------------------------%

\chapter{Rincian Kasus Uji}

\blindtext%


    %--------------------------------------------------------------------%
%
% Riwayat Hidup (Pedoman UM Lampiran 13) — naratif, spasi tunggal,
% tanpa foto. Batasi pendidikan sampai jenjang di atas SMA.
%
%--------------------------------------------------------------------%

\clearpage
\addcontentsline{toc}{chapter}{Riwayat Hidup}

\begin{center}
    \bfseries RIWAYAT HIDUP
\end{center}

\begin{singlespace}
\blindtext%
\end{singlespace}
     % Riwayat Hidup

\end{document}
